\documentclass[11pt,letterpaper]{article}

\usepackage[margin=0.78in,headheight=15pt]{geometry}
\usepackage[T1]{fontenc}
\usepackage{lmodern}
\usepackage{microtype}
\usepackage[table]{xcolor}
\usepackage{booktabs}
\usepackage{amsmath}
\usepackage{longtable}
\usepackage{tabularx}
\usepackage{array}
\usepackage{enumitem}
\usepackage{xurl}
\usepackage[colorlinks=true,linkcolor=olimexblue,urlcolor=olimexblue,citecolor=olimexblue]{hyperref}
\usepackage{fancyhdr}
\usepackage{listings}

\definecolor{olimexblue}{HTML}{1769AA}
\definecolor{olimexdark}{HTML}{263746}
\definecolor{lightblue}{HTML}{EAF3F9}
\definecolor{lightgray}{HTML}{F3F5F6}
\definecolor{warnbg}{HTML}{FFF4D6}
\definecolor{warnline}{HTML}{D49400}

\setlength{\parindent}{0pt}
\setlength{\parskip}{0.55em}
\setlist[itemize]{leftmargin=1.4em,itemsep=0.25em,topsep=0.25em}
\renewcommand{\arraystretch}{1.18}
\sloppy
\renewcommand{\thesection}{\arabic{section}.}
\renewcommand{\thesubsection}{\arabic{section}.\arabic{subsection}.}
\renewcommand{\thesubsubsection}{\arabic{section}.\arabic{subsection}.\arabic{subsubsection}.}
\setlength\LTleft{0pt}
\setlength\LTright{\fill}

\pagestyle{fancy}
\fancyhf{}
\cfoot{\thepage}
\renewcommand{\headrulewidth}{0pt}
\fancypagestyle{plain}{%
  \fancyhf{}
  \cfoot{\thepage}
  \renewcommand{\headrulewidth}{0pt}
}

\newcommand{\hexbyte}[1]{\texttt{0x#1}}
\newcommand{\field}[1]{\texttt{#1}}
\newcommand{\notice}[1]{%
  \begin{center}
  \fcolorbox{warnline}{warnbg}{\parbox{0.92\linewidth}{#1}}
  \end{center}
}
\newcolumntype{L}[1]{>{\raggedright\arraybackslash}p{#1}}
\newcolumntype{Y}{>{\raggedright\arraybackslash}X}

\lstset{
  basicstyle=\ttfamily\small,
  backgroundcolor=\color{lightgray},
  frame=single,
  rulecolor=\color{lightgray},
  columns=fullflexible,
  breaklines=true,
  showstringspaces=false,
  xleftmargin=0.4em,
  xrightmargin=0.4em
}

\hypersetup{
  pdftitle={OLinuXino EEPROM content - October 2026},
  pdfauthor={Olimex Ltd.},
  pdfsubject={Binary format for the OLinuXino A20 board identification EEPROM},
  pdfkeywords={OLinuXino, A20, EEPROM, board ID, CRC32}
}

\begin{document}

\begin{titlepage}
  \centering
  \vspace*{1.0in}
  {\Huge\bfseries\textcolor{olimexdark}{OLinuXino EEPROM content}\par}
  \vspace{0.24in}
  {\Large\textcolor{olimexblue}{A20 and T2 board-information data format}\par}
  \vspace{0.55in}
  {\Large Olimex Ltd.\par}
  \vspace{0.18in}
  {\large October 19, 2026\par}

  \vfill

  \begin{minipage}{0.84\textwidth}
  \centering
  \textbf{Status: Maintained replacement specification}\\[0.6em]
  This document corrects and supersedes the July 2018 and October 2019
  A20 EEPROM documents. It describes the 256-byte board-information record
  used by Olimex A20 and T2 OLinuXino products with a compatible EEPROM.
  \end{minipage}

  \vfill

  {\small
  Implementation baseline: Olimex U-Boot branch \field{release-20250305}\\
  The U-Boot source remains authoritative for the current board-ID registry.
  }
\end{titlepage}

\tableofcontents
\clearpage

\section{Scope}

This specification defines the first 256 bytes of the board-information EEPROM
used by supported Olimex boards based on the Allwinner A20 or T2 processor. It
covers the A20-OLinuXino-LIME, LIME2, MICRO, A20-SOM, A20-SOM204 and their T2
industrial variants when the hardware contains a compatible EEPROM.

The board module owns the identification record. A carrier or evaluation board
does not create a second board identity. For example, an A20-SOM204 mounted on
an evaluation board is identified by the EEPROM on the SOM204.

\notice{\textbf{Compatibility note.} Older hardware may have no EEPROM or may
contain erased or invalid data. Software must validate both the magic value and
CRC before using any field. A board name must be resolved from the current
U-Boot board-ID registry, not guessed from configuration bytes alone.}

\subsection{Physical EEPROM addressing}

The commonly used AT24C16-compatible EEPROM occupies I2C addresses
\hexbyte{50} through \hexbyte{57}. The data structure in this document occupies
the first 256-byte page, beginning at address \hexbyte{50}. Designers must
consider all eight addresses when adding devices to the same I2C/TWI bus.

\subsection{Normative implementation references}

The implementation corresponding to this revision is found in the Olimex
U-Boot repository:

\begin{itemize}
  \item \field{include/olimex/sys\_eeprom.h} - packed 256-byte structure.
  \item \field{board/olimex/common/sys\_eeprom.c} - EEPROM I/O and CRC validation.
  \item \field{include/olimex/boards.h} - storage, grade and capacity encodings.
  \item \field{board/olimex/common/boards.c} - current board-ID registry.
\end{itemize}

Repository: \url{https://github.com/OLIMEX/u-boot-olinuxino}

\section{General encoding rules}

\begin{itemize}
  \item The record length is exactly 256 bytes.
  \item All multi-byte integers are unsigned and stored little-endian.
  \item Text fields contain ASCII bytes and are not NUL-terminated unless
        explicitly stated.
  \item Capacity fields contain a base-2 exponent. A value $n$ represents
        $2^n$ bytes.
  \item Erased EEPROM bytes normally read as \hexbyte{FF}.
  \item Readers must reject a record with an invalid magic value or CRC.
\end{itemize}

\section{Memory layout}

\rowcolors{2}{lightblue}{white}
\begin{longtable}{L{0.85in} L{0.78in} L{1.45in} L{3.15in}}
\toprule
\textbf{Offsets} & \textbf{Length} & \textbf{Field} & \textbf{Encoding} \\
\midrule
\endhead
\hexbyte{00}-\hexbyte{03} & 4 & Header & Magic bytes
\field{55 AA 4C 4F} \\
\hexbyte{04}-\hexbyte{07} & 4 & Board ID & Unsigned 32-bit little-endian integer \\
\hexbyte{08}-\hexbyte{09} & 2 & Revision & Major and minor ASCII bytes \\
\hexbyte{0A}-\hexbyte{0D} & 4 & Serial number & Unsigned 32-bit little-endian integer \\
\hexbyte{0E}-\hexbyte{11} & 4 & Configuration & Storage, size, RAM and grade bytes \\
\hexbyte{12}-\hexbyte{1D} & 12 & Ethernet MAC & Twelve ASCII hexadecimal characters \\
\hexbyte{1E}-\hexbyte{FB} & 222 & Reserved & Reserved for future use \\
\hexbyte{FC}-\hexbyte{FF} & 4 & CRC-32 & Unsigned 32-bit little-endian CRC \\
\bottomrule
\end{longtable}
\rowcolors{2}{}{}

The MAC field ends at \hexbyte{1D}. Offset \hexbyte{1E} is the first reserved
byte.

\section{Field definitions}

\subsection{Header}

The first four bytes are fixed:

\noindent
\begin{tabular}{cccc}
\toprule
Offset & \hexbyte{00} & \hexbyte{01} & \hexbyte{02} \quad \hexbyte{03} \\
\midrule
Value & \hexbyte{55} & \hexbyte{AA} & \hexbyte{4C} \quad \hexbyte{4F} \\
\bottomrule
\end{tabular}

Interpreted as a little-endian 32-bit integer, these bytes equal
\hexbyte{4F4CAA55}. A different value means the record is invalid.

\subsection{Board ID}

The board ID at offsets \hexbyte{04}-\hexbyte{07} is a unique unsigned 32-bit
identifier stored little-endian. It selects the product name, expected storage,
RAM, grade and device-tree configuration.

Board IDs are maintained in \field{board/olimex/common/boards.c}. Applications
should update their lookup table from that source. Unknown IDs must be reported
as unknown while the remaining structurally valid fields may still be shown.

\clearpage
\subsection{Revision}

\noindent
\rowcolors{2}{lightblue}{white}
\begin{tabular}{lll}
\toprule
Offset & Meaning & Valid value \\
\midrule
\hexbyte{08} & Major revision & ASCII \field{A} through \field{Z} \\
\hexbyte{09} & Minor revision & ASCII \field{1} through \field{9}, or unused \\
\bottomrule
\end{tabular}
\rowcolors{2}{}{}

U-Boot writes \hexbyte{00} when the minor revision is unused. Some production
or erased records contain \hexbyte{FF}. Readers should display no minor
revision for any byte outside ASCII \field{1} through \field{9}.

Examples are \field{C}, \field{G2} and \field{J1}.

\subsection{Serial number}

Offsets \hexbyte{0A}-\hexbyte{0D} contain a sequential unsigned 32-bit serial
number in little-endian order. U-Boot commonly displays it as eight hexadecimal
digits.

\subsection{Configuration}

The configuration occupies four independent bytes:

\noindent
\rowcolors{2}{lightblue}{white}
\begin{tabular}{lll}
\toprule
Offset & Field & Meaning \\
\midrule
\hexbyte{0E} & Storage & Primary onboard nonvolatile storage type \\
\hexbyte{0F} & Storage size & Capacity exponent, valid when storage is not \hexbyte{00} \\
\hexbyte{10} & RAM size & Capacity exponent \\
\hexbyte{11} & Grade & Temperature-grade classification \\
\bottomrule
\end{tabular}
\rowcolors{2}{}{}

\subsubsection{Storage type}

\noindent
\rowcolors{2}{lightblue}{white}
\begin{tabular}{lll}
\toprule
Value & ASCII & Meaning \\
\midrule
\hexbyte{00} & - & No onboard storage represented by this field \\
\hexbyte{65} & \field{e} & eMMC \\
\hexbyte{6E} & \field{n} & NAND flash \\
\hexbyte{73} & \field{s} & SPI flash \\
\bottomrule
\end{tabular}
\rowcolors{2}{}{}

The letters are lowercase. Uppercase \field{E}, \field{N} and \field{S} are
not the encoding used by current Olimex U-Boot.

Some product names describe more than one fitted storage device, for example
\field{e16Gs16M}. The configuration byte records the primary storage type;
U-Boot identifies additional SPI flash from the board ID.

\clearpage
\subsubsection{Capacity encoding}

For both storage size and RAM size:

\[
  \text{capacity in bytes} = 2^{\text{field value}}
\]

\noindent
\rowcolors{2}{lightblue}{white}
\begin{tabular}{llll}
\toprule
Value & Capacity & Value & Capacity \\
\midrule
0 & 1 B & 10 & 1 KiB \\
20 & 1 MiB & 24 & 16 MiB \\
29 & 512 MiB & 30 & 1 GiB \\
32 & 4 GiB & 33 & 8 GiB \\
34 & 16 GiB & & \\
\bottomrule
\end{tabular}
\rowcolors{2}{}{}

Product names traditionally use MB and GB even though the encoding is based on
powers of two.

\subsubsection{Grade}

\noindent
\rowcolors{2}{lightblue}{white}
\begin{tabular}{ll}
\toprule
Value & Meaning \\
\midrule
\hexbyte{00} & Commercial grade \\
\hexbyte{01} & Industrial grade \\
\bottomrule
\end{tabular}
\rowcolors{2}{}{}

\subsection{Ethernet MAC address}

Offsets \hexbyte{12}-\hexbyte{1D} contain exactly twelve uppercase ASCII
hexadecimal characters, without separators. For example:

\noindent
\field{301F9AD0ACB6} represents \field{30:1F:9A:D0:AC:B6}.

If a valid MAC address was not programmed, the field can contain twelve
\field{F} characters or erased \hexbyte{FF} bytes. Readers should display
\field{FF:FF:FF:FF:FF:FF} and must not install it as an Ethernet address.

Original A20-SOM products do not contain an onboard Ethernet PHY and normally
do not carry a usable Ethernet MAC. A20-SOM204 behavior is variant-dependent.
In the current U-Boot registry, IDs 8958, 10257 and 10234 are classified as
having EMAC. Software should use the current board-ID registry rather than
assuming that every SOM204 variant has or lacks Ethernet.

\subsection{Reserved area}

Offsets \hexbyte{1E}-\hexbyte{FB} are reserved. Current U-Boot initializes this
area to \hexbyte{FF}. Readers must ignore it. Writers should initialize it to
\hexbyte{FF} for a new record and should not assign private meanings that could
conflict with future revisions.

\subsection{CRC-32}

Offsets \hexbyte{FC}-\hexbyte{FF} contain the CRC-32 of bytes
\hexbyte{00}-\hexbyte{FB}, stored as a little-endian unsigned 32-bit value.

The algorithm is the U-Boot \field{crc32()} implementation invoked as:

\begin{lstlisting}[language=C]
crc = crc32(0L, record, 252);
\end{lstlisting}

The CRC field itself is excluded from the calculation. A record is valid only
when the stored value equals the calculated value.

\section{Validation procedure}

A reader should perform these steps in order:

\begin{enumerate}[leftmargin=1.6em]
  \item Read the full 256-byte page beginning at I2C address \hexbyte{50}.
  \item Confirm bytes 0-3 are \field{55 AA 4C 4F}.
  \item Calculate CRC-32 over bytes 0-251 with initial value zero.
  \item Decode the stored CRC from bytes 252-255 as little-endian.
  \item Reject the record if the CRC values differ.
  \item Decode the remaining fields.
  \item Resolve the board ID using the current board registry.
\end{enumerate}

A diagnostic tool may offer an explicit option to decode a record with a bad or
unavailable CRC, but it must clearly mark the result as unverified.

\section{Worked configuration examples}

\noindent
\rowcolors{2}{lightblue}{white}
\begin{tabularx}{\textwidth}{L{1.55in} L{0.9in} L{0.9in} L{0.9in} Y}
\toprule
Product configuration & Storage byte & Size byte & RAM byte & Grade byte \\
\midrule
No storage, 1 GiB RAM, commercial &
\hexbyte{00} & Not applicable & 30 (\hexbyte{1E}) & \hexbyte{00} \\
4 GiB eMMC, 1 GiB RAM, commercial &
\hexbyte{65} & 32 (\hexbyte{20}) & 30 (\hexbyte{1E}) & \hexbyte{00} \\
16 GiB eMMC, 1 GiB RAM, commercial &
\hexbyte{65} & 34 (\hexbyte{22}) & 30 (\hexbyte{1E}) & \hexbyte{00} \\
16 MiB SPI, 512 MiB RAM, industrial &
\hexbyte{73} & 24 (\hexbyte{18}) & 29 (\hexbyte{1D}) & \hexbyte{01} \\
\bottomrule
\end{tabularx}
\rowcolors{2}{}{}

\clearpage
\section{A20-SOM and A20-SOM204 board IDs}

This table is an informative snapshot from Olimex U-Boot
\field{release-20250305}. The source registry is authoritative.

\small
\rowcolors{2}{lightblue}{white}
\begin{longtable}{@{}L{0.65in}@{\hspace{0.16in}}L{3.40in}@{\hspace{0.16in}}L{1.20in}@{\hspace{0.16in}}L{1.19in}@{}}
\toprule
\textbf{ID} & \textbf{Board name} & \textbf{Storage} & \textbf{Grade} \\
\midrule
\endhead
4673 & A20-SOM-n4GB & NAND 4 GiB & Commercial \\
7664 & A20-SOM & None & Commercial \\
8849 & A20-SOM-IND & None & Industrial \\
8922 & A20-SOM-n8GB & NAND 8 GiB & Commercial \\
9155 & A20-SOM-e16GB & eMMC 16 GiB & Commercial \\
9148 & A20-SOM-e16GB-IND & eMMC 16 GiB & Industrial \\
9047 & A20-SOM-e16Gs16M & eMMC 16 GiB & Commercial \\
9259 & T2-SOM-IND & None & Industrial \\
9827 & T2-SOM-e8Gs16M-IND & eMMC 8 GiB & Industrial \\
11454 & T2-SOM-e16Gs16M-IND & eMMC 16 GiB & Industrial \\
\midrule
8991 & A20-SOM204-1G & None & Commercial \\
8958 & A20-SOM204-1Gs16Me16G-MC & eMMC 16 GiB & Commercial \\
10257 & A20-SOM204-1G-M & None & Commercial \\
10157 & T2-SOM204-1Gs16Me4G-C-I & eMMC 4 GiB & Industrial \\
10234 & T2-SOM204-1Gs16Me8G-MC-I & eMMC 8 GiB & Industrial \\
11458 & T2-SOM204-1Gs16Me16G-M-I & eMMC 16 GiB & Industrial \\
11462 & T2-SOM204-1Gs16Me16G-MC-I & eMMC 16 GiB & Industrial \\
10238 & T2-SOM204-1G-I & None & Industrial \\
\bottomrule
\end{longtable}
\rowcolors{2}{}{}
\normalsize

\section{Change-control guidance}

Future changes should commit both this LaTeX source and its generated PDF.
Changes to the binary layout require a corresponding implementation change and
a new document release date. Updates to product IDs normally belong only in
\field{boards.c}; this document should include a dated snapshot only when it is
useful for users.

\end{document}
